%======================================================================
% CAPÍTULO 15 - ALGORITMO DE FUSIÓN (HASH JOIN)
%======================================================================
\chapter{Algoritmo de Fusión (Hash Join)}

\bigskip

El algoritmo de fusión es responsable de unir las tres fuentes de datos (Base, Movimientos, Beneficiarios) en un expediente digital unificado. Este capítulo detalla la implementación del Hash Join en memoria.

\bigskip

\section{Descripción del Problema}

El sistema debe unir tres fuentes de datos distintas:

\bigskip

\begin{enumerate}
    \item Beneficiarios: Stream principal de registros.
    \item Base: Datos estructurales del personal.
    \item Movimientos: Datos financieros transaccionales.
\end{enumerate}

\bigskip

El desafío principal es realizar esta fusión de manera eficiente para 500,000+ registros. Las consultas SQL tradicionales (JOIN) resultan extremadamente costosas en términos de tiempo cuando se procesan grandes volúmenes.

\bigskip

\section{Implementación del Hash Join}

Sentinel implementa una variante del algoritmo Hash Join en memoria, aprovechando la velocidad de acceso a RAM para lograr tiempos de procesamiento del orden de segundos.

\bigskip

\subsection{Fase de Indexación (Build Phase)}

\bigskip

En la fase de indexación, se cargan las entidades Base y Movimientos en memoria y se construyen tablas hash optimizadas. La complejidad de esta fase es:

\[
\mathcal{O}(n + m)
\]

donde n es el número de registros Base y m es el número de registros de Movimientos.

\bigskip

\begin{enumerate}
    \item Se cargan los registros Base desde gRPC.
    \item Se crea un HashMap para búsqueda O(1) con la clave patterns.
    \item Se cargan los registros de Movimientos.
    \item Se crea un HashMap con la clave cédula.
\end{enumerate}

\bigskip

Ejemplo de código de indexación:
\begin{verbatim}
let mut base_index: HashMap<String, &Base> = HashMap::new();
for base in bases.iter() {
    base_index.insert(base.patterns.clone(), &base);
}

let mut movimiento_index: HashMap<String, Vec<&Movimiento>> = HashMap::new();
for movimiento in movimientos.iter() {
    movimiento_index
        .entry(movimiento.cedula.clone())
        .or_insert_with(Vec::new)
        .push(&movimiento);
}
\end{verbatim}

\bigskip

\subsection{Fase de Sondeo (Probe Phase)}

\bigskip

En la fase de sondeo, el stream de Beneficiarios entra al sistema y para cada registro se realiza una búsqueda O(1) en los índices. La complejidad total es:

\[
\mathcal{O}(k)
\]

donde k es el número de Beneficiarios.

\bigskip

\begin{enumerate}
    \item Se recibe el stream de Beneficiarios.
    \item Para cada Beneficiario, se busca su Base por patterns.
    \item Se buscan sus Movimientos por cédula.
    \item Se fusionan en un Beneficiario enriquecido.
\end{enumerate}

\bigskip

Ejemplo de código de sondeo:
\begin{verbatim}
for beneficiario in beneficiarios {
    // Búsqueda O(1) en índice Base
    if let Some(base) = base_index.get(&beneficiario.patterns) {
        // Búsqueda O(1) en índice Movimientos
        if let Some(movs) = movimiento_index.get(&beneficiario.cedula) {
            // Fusionar
            let fusionado = BeneficiarioEnriquecido {
                base: base.clone(),
                movimientos: movs.clone(),
                ..beneficiario
            };
            resultados.push(fusionado);
        }
    }
}
\end{verbatim}

\bigskip

\section{Complejidad Algorítmica}

\begin{center}
\begin{tabular}{lll}
\toprule
\textbf{Operación} & \textbf{Complejidad} & \textbf{Descripción} \\
\midrule
Indexación Base & O(n) & Crear HashMap de Base \\
Indexación Movimientos & O(m) & Crear HashMap de Movimientos \\
Sondeo & O(k) & Búsqueda O(1) por registro \\
\midrule
\textbf{Total} & \textbf{O(n + m + k)} & Lineal \\
\bottomrule
\end{tabular}
\end{center}

\bigskip

\section{Teorema de Optimalidad}

El algoritmo Hash Join en memoria es óptimo para el problema de fusión de entidades cuando:

\begin{enumerate}
    \item Los datos caben en memoria.
    \item La clave de búsqueda tiene buena distribución (baja colisión).
    \item Se requiere acceso aleatorio a los datos.
\end{enumerate}

\bigskip

\section{Ventajas del Enfoque}

El enfoque In-Memory Hash Join ofrece las siguientes ventajas sobre los JOINs tradicionales:

\bigskip

\begin{itemize}
    \item Eliminación de latencia de I/O: Sin consultas a base de datos durante el procesamiento.
    \item Búsqueda O(1): Tiempo de acceso constante independiente del tamaño de los datos.
    \item Paralelizable: Las búsquedas pueden distribuirse entre múltiples núcleos.
    \item Reproducible: Los mismos datos siempre producen el mismo resultado.
\end{itemize}

\vfill
\newpage
